% Side-by-side reasoning-trace example for the methods classifier:
% v1 schema (methods_01.csv) vs v3 schema (methods_04.csv) on the same paper.
% Needs \usepackage[most]{tcolorbox} and \usepackage{hyperref}.
% Reasoning text is verbatim from the model snapshots, lightly abridged.
%
% Model: Gemma-4-31B IT (thinking mode, temperature 0, vLLM serving).
% v1 categories: qual, quant, both, Descriptive, spatial/mapping, unclear.
% v3 categories: qual, quant, mixed, spatial, descriptive-empirical,
%                theoretical-conceptual, unclear.
% v1 final: spatial/mapping. v3 final: spatial. Annotator: quant.
% Source runs: 2026-05-13_1512_methods_v1_abstract,
%              2026-05-13_2233_methods_v3_abstract_cat04.

\begin{figure*}[t]
  \centering

  \textbf{Paper.} ``Geographical Disparities in Energy Access: Challenges and Opportunities for Ghana's Energy Transition''
  (\href{https://doi.org/10.1080/00330124.2025.2582789}{doi:10.1080/00330124.2025.2582789}).

  \smallskip
  \textbf{Methodology sentence in abstract.} \emph{``Using a mixed-methods approach combining geospatial analysis with statistical evaluation of regional data.''}

  \smallskip
  \textbf{Annotator label.} \texttt{quant}.

  \medskip

  \begin{tcbraster}[
    raster columns=2,
    raster equal height=rows,
    raster column skip=8pt,
    raster row skip=4pt,
  ]
  % ----------------------------- v1 box --------------------------------------
  \begin{tcolorbox}[
    enhanced,
    colback=red!2,
    colframe=red!50!black!55,
    boxrule=0.4pt, arc=2pt,
    left=8pt, right=8pt, top=6pt, bottom=6pt,
    title={\textsc{v1 schema} (methods\_01.csv) --- 6 categories},
    fonttitle=\small\bfseries,
    fontupper=\small,
  ]
  \begin{itemize}\setlength\itemsep{2pt}
  \item \textbf{Reading \texttt{both}.} ``\emph{`Mixed-methods'} usually refers to Qual~+~Quant. Here, it says \emph{`geospatial analysis'}~+~\emph{`statistical evaluation'}.''
  \item \textbf{Is this \texttt{both}?} ``No, because \emph{`statistical evaluation'} is quantitative. So it's Quant~+~Spatial.''
  \item \textbf{Fits \texttt{spatial/mapping}?} ``Yes, it uses geospatial analysis to analyze spatial patterns and proximity.''
  \item \textbf{Fits \texttt{quant}?} ``Yes, it uses statistical evaluation.''
  \item \textbf{Reconciling.} ``The guidelines say: \emph{`Use spatial/mapping when the analytical method is fundamentally spatial: GIS, remote sensing, spatial analysis\dots{} OR a spatial statistical method'.}''
  \item \textbf{Decision rule.} ``If a paper combines standard stats with GIS/spatial analysis, and the core focus is on `geographical disparities' and `spatial patterns', it leans heavily towards \texttt{spatial/mapping}.''
  \item \textbf{Re-checking \texttt{both}.} ``The \texttt{both} definition is \emph{`uses both qualitative and quantitative approaches substantively'}. This paper does \emph{not} use qualitative methods.''
  \end{itemize}

  \medskip
  \textbf{Final prediction.}\quad\texttt{spatial/mapping}
  \end{tcolorbox}

  % ----------------------------- v3 box --------------------------------------
  \begin{tcolorbox}[
    enhanced,
    colback=green!2,
    colframe=green!40!black!55,
    boxrule=0.4pt, arc=2pt,
    left=8pt, right=8pt, top=6pt, bottom=6pt,
    title={\textsc{v3 schema} (methods\_04.csv) --- 7 categories, sharpened},
    fonttitle=\small\bfseries,
    fontupper=\small,
  ]
  \begin{itemize}\setlength\itemsep{2pt}
  \item \textbf{Reading \texttt{mixed}.} ``The abstract says \emph{`mixed-methods approach'}, but it defines this as \emph{`combining geospatial analysis with statistical evaluation'}. In the provided labels, \texttt{mixed} is defined as \emph{`qualitative interviews/ethnography AND quantitative analysis'}. Since there is no qualitative work, \texttt{mixed} in the academic sense of Qual~+~Quant is not applicable here.''
  \item \textbf{Ruling out \texttt{qual}.} ``No mention of interviews, ethnography, or focus groups.''
  \item \textbf{Identifying spatial signal.} ``The abstract mentions \emph{`geospatial analysis'}, \emph{`spatial patterns'}, \emph{`spatial inequities'}, and \emph{`proximity to power generation facilities'}.''
  \item \textbf{\texttt{quant} vs \texttt{spatial} criteria.} ``\texttt{quant} is for \emph{`statistical analysis on a large-N survey or non-spatial numeric dataset --- no spatial methods'}. \texttt{spatial} is for \emph{`explicit spatial-statistical method\dots{} or a GIS overlay that produces the inferential claim'}.''
  \item \textbf{Decision.} ``The abstract says \emph{`combining geospatial analysis with statistical evaluation'}. The focus is on \emph{`spatial patterns'} and \emph{`geographical disparities'}. This fits \texttt{spatial} better than \texttt{quant} because the spatial component is central to the analysis, not just a side note.''
  \end{itemize}

  \medskip
  \textbf{Final prediction.}\quad\texttt{spatial}
  \end{tcolorbox}
  \end{tcbraster}

  \caption{Reasoning traces from Gemma-4-31B (thinking mode, temperature
    $=0$) on the same borderline paper, under two schemas. The v1 schema's
    \texttt{both} category --- ``mixed methods'' phrased as ``uses both
    qualitative and quantitative approaches'' --- forces the model into a
    long deliberation: it must rule out the \texttt{both} interpretation,
    reconsider it, and then choose between \texttt{spatial/mapping} and
    \texttt{quant} by re-reading the \texttt{spatial/mapping} definition.
    The v3 schema replaces \texttt{both} with \texttt{mixed} and pins it to
    \emph{``qualitative interviews/ethnography AND quantitative analysis''},
    which lets the model dismiss the qualitative interpretation in a single
    step and proceed directly to the \texttt{spatial} vs \texttt{quant}
    distinction. Both schemas yield the same final prediction (\texttt{spatial})
    on this paper; the annotator chose \texttt{quant}. The illustrative
    point is the reasoning structure, not the agreement: sharpened
    definitions produce a shorter, more legible deliberation even when the
    final classification is unchanged.}
  \label{fig:reasoning-trace-ghana}
\end{figure*}
